\section{The common global ledger and its final uniform installation}
\label{sec:global-final-installation-v035}

\status{This chapter constructs the actual tied ledger from the nine
original numerical and lattice predicates.  Together with LFS it
completes the remaining global assembly implication.  The common
family, compatible second layer, duties, and arithmetic and spectral
producer inputs retain their original scope.  No new family or
unrelated model is chosen to satisfy a missing condition.}

Throughout, $\mathcal F:\texttt{FamP}(\epsilon)$ is the supplied
original family, with $0<\epsilon<1/52$.  At a scale $x>\mathcal F.x_1$
write $D_i=\mathcal F.D(x,i)$ and $L_i=\mathcal F.L(x,i)$.
We distinguish the original record $L_i:\texttt{SizeData}_2(D_i)$
from the eleven-predicate \texttt{LChain.SizeData}$_2$ on a numerical
ledger.  The latter is the input to ZW1--ZW4.  All eventual thresholds
below are chosen before the member, base point, word, or finite place.

% BEGIN GFI1
\paragraph{GFI1: the common family dictionary includes arithmetic data.}
\label{claim:GFI1}
The original fields \texttt{R\_eq}, \texttt{L\_x},
\texttt{L\_eps} and \texttt{L\_J} imply that all $D_i$ have the
same ranks and frame, and all $L_i$ have the same scale.  Indeed, a
scale is specified by $x,\epsilon,J$; its remaining components are
functions of those values.  Hence $a,b,b',k,n,m,f,L$ agree, where
$L=x^\epsilon$, $b=e^L$, $a=xb$ and $f=2n$.
The separate original family field \texttt{eta\_eq} gives the
common value of $L_i.\mathrm{etaSz}$; it is not inferred from
the scale alone.  This is the equality used by \texttt{TieGlobal}
and the remaining-size comparison in GFI12.

The map from the original top set to integer labels is injective and
has image $T_n=\mathcal F.TTn(x)$.  It therefore gives a bijection
\[
 D_i.\mathrm{host}.TT\simeq T_n,
 \qquad \#D_i.\mathrm{host}.TT=\#T_n.                 \tag{GFI.1}
\]
The second-set integer labels agree by
\texttt{realisedChain\_torus\_eq}.  The family also fixes $q,d,QS$,
the root discriminant, $g$, and $C_\Gamma$.  Consequently it fixes
$v$, $\log(g_0g_1)$, $L_a$, $B_0$, $N_c$ and $\Gamma$.
For $N_c$ use the actual second-set equality before taking its
cardinality.  For $\Gamma$ use
$\exp(C_\Gamma m^5 fL)$.  The quadratic logarithmic derivative
is fixed because its character is defined by the same prime $d$;
this identifies the actual value at one, including its real part.

There is also equality of the arithmetic records after transporting
along the frame equality.  Their modulus is the same $QSqd$.
The field \texttt{art\_eq} gives the same signed reciprocity value
and cyclic integer label at every integer coprime to that modulus.
Every residue unit is represented by its least nonnegative integer,
which is coprime to the modulus.  Thus the two reciprocity
homomorphisms agree on every unit.  Injectivity of the cyclic labels
and of the two sign codes recovers equality of the actual values.
The definition of \texttt{arithOf} then gives equality of the
arithmetic records, with the prescribed interchange of the two sign
coordinates.  This is equality on the source objects, rather than a
choice of a second reciprocity model.
% END GFI1

% BEGIN GFI2
\paragraph{GFI2: carry one realized degree pair to every base point.}
\label{claim:GFI2}
Suppose the original \texttt{LossFS} conclusion is available at $x$.
It supplies $i_0$, integer functions $a_v,b_v$, their admissible
windows, and an index $i_1$ realizing those functions on every top.
It also supplies \texttt{Loss\_at} for every member realizing them.
Let $H=D_{i_0}.R.A$.  For each $h\in H$, its integer label is below
$n$.  The family's \texttt{coverS} property applied to $i_1$ gives
an index $\mathrm{pt}(h)$ whose distinguished base point has that
label and whose top and bottom degrees agree with those at $i_1$
whenever the top labels agree.

Every top at $\mathrm{pt}(h)$ has a counterpart at $i_1$, by
(GFI.1).  Their degrees are therefore exactly
$a_v(\mathrm{val}(t))$ and $b_v(\mathrm{val}(t))$.  Choose these
indices for the finite set $H$.  We have constructed the literal
\texttt{Mod} record
\[
 M=(i_0,a_v,b_v,\mathrm{pt}),                            \tag{GFI.2}
\]
including both its base-point equality and its degree-carrying
condition.  All members $\mathrm{pt}(h)$ consequently satisfy the
same loss estimate.  Only the actual realized pair is transported;
no assertion that every admissible pair is realized is used.
% END GFI2

% BEGIN GFI3
\paragraph{GFI3: top sets, place labels, and weights on that model.}
\label{claim:GFI3}
Use the original frame $e=D_{i_0}.R.\mathrm{toFrame}$ and put
\[
 T=\{h\in H:\mathrm{val}(h)\in T_n\},\qquad
 T_2=e.\mathrm{dig}(D_{i_0}.\mathrm{host}.T2),\qquad
 a_h=a_v(\mathrm{val}(h)).                              \tag{GFI.3}
\]
The inverse of the integer-label map gives an equivalence
$e_T:T_n\simeq T$.  Hence $\#T=\#T_n$.  The original second-set
entries lie below the digit map's injectivity bound: its stored
range condition is $4L_w s<n_0$, and $L_w\ge2$.  Thus
\[
 \#T_2=\#D_{i_0}.\mathrm{host}.T2.                     \tag{GFI.4}
\]
The same set describes the second-set image for every $D_{\mathrm{pt}(h)}$
by GFI1.

The original place coding maps its $\tau$-set bijectively to
$\Phi\subset E_0\times H$.  Injectivity follows from injectivity
of both sign codes, the cyclic labels, and the original $e_\tau$;
surjectivity follows from membership and the common cardinality
$\#\Phi=2n=f$.  Choose the source's finite enumeration
$e_P:\{0,\ldots,f-1\}\simeq\Phi$.

For $h\in H$, $0\le j<m$ and $\rho\in E_0\times H$, define
$w_{h,j}(\rho)$ to be the original $w_{j+1}(\tau)$ if $\rho$
is the code of a place of $D_{\mathrm{pt}(h)}$, and to be $1$
otherwise.  The representing $\tau$ is unique, so this is the
literal source weight.  Reindexing gives, for any function $F$,
\[
 \sum_{r=0}^{f-1}F(w_{h,j}(e_P(r)))
 =\sum_\tau F(w_{j+1}(\tau)).                           \tag{GFI.5}
\]
Similarly the lifted indices are exactly $1,\ldots,m$, and
\[
 \sum_{t\in D_{\mathrm{pt}(h)}.\mathrm{host}.TT}F(a_T(t))
 =\sum_{t\in T_n}F(a_v(t)).                              \tag{GFI.6}
\]
These are bijective reindexings of the actual sums.  They preserve
their multiplicities and all metric-dependent quantities.
% END GFI3

% BEGIN GFI4
\paragraph{GFI4: construct every entry of the original numerical ledger.}
\label{claim:GFI4}
Let $W_h$ be the original word set of $D_{\mathrm{pt}(h)}$.
It is finite and nonempty, with its stored distinguished word.
For a cutoff $N$, constants $C_H,\bar\eta,\bar e$, and error
functions $\eta_h,e_h$, define the ledger $\mathscr D$ as follows.
Its scalar entries $x,\epsilon,a,b,b',k,L,L_a,B_0,N_c,v,
\log(g_0g_1),\log\Gamma$ are those at $i_0$; its Euler constant is
$\gamma$ and its logarithmic derivative is $L_{i_0}$'s actual one.
Its top and place weights are (GFI.3) and (GFI.5).

For $u\in W_h$, put
\[
\begin{aligned}
 d_{h,u}&=\operatorname{rank}(D_{\mathrm{pt}(h)}.\mathrm{div.Nu}(u)),&
 T_{h,u}&=(D_{\mathrm{pt}(h)}.\mathrm{mem}(u)).A.T,\\
 \mu_{N,h,u}&=\mu(D_{\mathrm{pt}(h)}.\mathrm{div.Nu}(u)),&
 \mu^{\max}_{N,h,u}&=\mu_{\max}(D_{\mathrm{pt}(h)}.\mathrm{div.Nu}(u)),\\
 \mu_{M,h,u}&=\mu((D_{\mathrm{pt}(h)}.\mathrm{fm}(u)).\mathrm{lat}(0)),&
 \mu^{\max}_{M,h,u}&=\mu_{\max}((D_{\mathrm{pt}(h)}.\mathrm{fm}(u)).\mathrm{lat}(0)),\\
 h_{h,u}&=D_{\mathrm{pt}(h)}.\mathrm{div.h}(u),&
 E_{h,u,j}&=\sum_{w\in\mathrm{supp}}\mathrm{sz}(w,E(u,j+1,w)).
\end{aligned}                                                     \tag{GFI.7}
\]
The last sum uses $L_{\mathrm{pt}(h)}$'s original support and
valuation normalization.  Set $\mu_{C,h}=\mu(RR^\vee)$ and
$\mu^{\min}_{C,h}=\mu_{\min}(RR^\vee)$ for that same member.
Copy its $\mathrm{Fo}$ and $\mathrm{Ga}$ and set
\[
 \mathrm{Ta}_h=\sum_{w\in\mathrm{supp}}\mathrm{sz}(w,\tau_w),
 \qquad
 \mathrm{IO}_h=\sum_{w\in\mathrm{supp}}\ell_we_w\,\mathrm{io}_w.
                                                               \tag{GFI.8}
\]
Thus $\mathrm{IO}$ has the source's single $\ell_we_w$ factor;
it is not confused with the doubled squared-metric sum.
Using the common arithmetic record $A$ from GFI1, set
\[
 \mathrm{ASK}_h=\sum_{t\in T}\sum_{j=0}^{m-1}
       \mathrm{ASK}(A,N,a_t,t,w_{h,j}),\qquad HH_h=x^{C_H}b.
                                                               \tag{GFI.9}
\]
Finally copy the actual retained defect $\delta_h$, put
$\bar\delta=e^{-x^{\epsilon/(7J)}}$, and use the chosen
$\eta_h,e_h,\bar\eta,\bar e$.

These formulas specify every field of the source's \texttt{Data}.
GFI1 proves \texttt{TieGlobal} at every $h$, including the common
$X''$, $\eta_{\mathrm{Sz}}$ and $X'$.  Formulae (GFI.7)--(GFI.8)
prove every \texttt{TieLocal} field by definition.  GFI2 and the
common ranks and primes prove \texttt{TiePts}.  The source and target
lattices, rational maps, metrics and word sets are retained without
any change of model.
% END GFI4

% BEGIN GFI5
\paragraph{GFI5: one character functional and an admissible cutoff.}
\label{claim:GFI5}
Take the original character functional $\ell$ supplied by GFA12
at $i_0$, for the common arithmetic record $A$.  It agrees with the
logarithmic derivative of every primitive character realizing each
nonzero additive character.  It also satisfies
$L_{i_0}.\mathrm{logDerivativeAtOne}=\operatorname{Re}\ell(\mathrm{chC})$.
Keep this single $\ell$ for all $h$.

GFI3, the unchanged second set, and that character identity prove
\texttt{TieArith}.  The equalities of top weights, place weights,
Euler constant, quadratic logarithmic derivative, and the ASK sum
in (GFI.9) are exactly the five fields of \texttt{Ident}.

The original node estimate of GFA10 gives $w_{h,j}(\rho)\ge1$
for $\rho\in\Phi$.  Each common top label can be read at the
single member $\mathrm{pt}(0)$; its original top window and GFA11
give
\[
 a\le a_t\le2a,\qquad |a_t/b-x|\le1.                \tag{GFI.10}
\]
Choose $N=\lceil2a\rceil_{\mathbb N}$.  Then
$a_t/w_{h,j}(\rho)\le a_t\le N$ for all permitted $h,j,t,\rho$.
Moreover the five original \texttt{Reads} conditions hold:
frame and rank equality are GFI1, top membership is (GFI.1), top
degree equality is GFI2, and place-weight equality is GFI3.
The arithmetic-record equality is the one proved in GFI1.  Thus
the loss statement can actually be applied to (GFI.9).
% END GFI5

% BEGIN GFI6
\paragraph{GFI6: a common threshold, constants, and positivity record.}
\label{claim:GFI6}
Assume the nine original family predicates
\[
\begin{gathered}
 \texttt{CountF},\ \texttt{SlopesMF},\ \texttt{SlopesNF},\
 \texttt{LatF},\ \texttt{NodesF},\ \texttt{FrameF}(C_H),\\
 \texttt{LossFS}(e_L),\ \texttt{CharF},\ \texttt{WinF},
\end{gathered}                                                     \tag{GFI.11}
\]
together with the original \texttt{ChosenTT},
\texttt{ChosenQD}$(C'_5)$, \texttt{ChosenThin}$(c_6)$,
\texttt{LinkIRR} and \texttt{LinkDIV}$(C_{\rm abs},C_H)$.
Assume $C_{\rm abs},C_H,e_L\ge0$ and $c_6>0$.
Extract their constants and thresholds first, using
\texttt{CharF} at $(c_6,C'_5)$ and \texttt{WinF} at $c_6$.
Write $C_c,C_s,C_1,C_\ell\ge0$ for the counting, two slope,
and loss constants, and $C_d$ for the divisibility constant.
Put
\[
 C_d^+=\max(C_d,0),\quad
 K=C_d^++C_c+C_s+C_1+C_\ell+6,\quad
 c_{HH}=C_H,\quad C_L=C'_5.                              \tag{GFI.12}
\]
The ledger constant record has, in source order,
\[
 (c_D,e_D,c_E,c_H,c_C,c_S,c_{\mathrm{sl}},c_P,e_P,c_{HH})
 =(C_d^+,C_{\rm abs},C_d^+,C_d^+,C_c,C_s,C_1,C_\ell,e_L,C_H).
                                                               \tag{GFI.13}
\]
Here $c_H$ in the record is a height-comparison coefficient,
distinct from the exponent $C_H$.

Take one threshold greater than $\mathcal F.x_1$ and all extracted
thresholds.  For example, $\mathcal F.x_1+1$ plus the sum of
their absolute values works.  Only then fix $x$ in the original
window, choose (GFI.2), and use \texttt{FrameF} to choose
$\bar\eta,\bar e$ before any member.  For every $h$ choose its
frame errors $\eta_h,e_h$ within these same bounds.  They satisfy
\[
 0\le\bar e\le\tfrac14,\quad 0\le\bar\eta,\qquad
 (\bar\eta+4\sqrt{\bar e})x^{C_H}\le6.                 \tag{GFI.14}
\]

The ledger now satisfies all seventeen positivity conditions in
\texttt{HidaSetting}.  Indeed $x,b,a,b',k,v,f$ have their
required signs; $a_t\ge0$ follows from (GFI.10), and place weights
are positive by GFI5.  Actual ranks are positive by \texttt{LatF},
and actual level primes are positive.  The same predicate gives
$\mu_N\le\mu_N^{\max}$.  Also $C_1\ge0$ and
$\bar\delta=e^{-x^{\epsilon/(7J)}}<1$, since $J\ge1$ and
$\epsilon>0$.  The remaining three error conditions are (GFI.14).
No positivity of an artificial rank or replacement form is used.
% END GFI6

% BEGIN GFI7
\paragraph{GFI7: divisibility and the two comparison fields.}
\label{claim:GFI7}
The last three estimates in the original \texttt{LinkDIV} record
give \texttt{DIV}, \texttt{EComp5}, and \texttt{HComp} for
$\mathscr D$.  Here is the complete transport and constant check.
First $L_a,B_0,N_c\ge0$ by the original windows and their
definitions.  All the quantities multiplying $C_d$ in those three
estimates are therefore nonnegative.  Replacing $C_d$ by $C_d^+$
is consequently valid even if the supplied $C_d$ was negative.

For \texttt{DIV}, GFI1 identifies the original
$\log\Omega=\gamma+\log(\pi/v)-L'/L(1)$ at every member.
Using (GFI.6), its main term becomes exactly
\[
 \sum_{t\in T_n}\bigl[m(\phi(a_v(t))-a_v(t)\log\Omega)
                       -a_v(t)\log T_{h,u}\bigr],
 \qquad \phi(y)=y\log y-y.                            \tag{GFI.15}
\]
The other terms are $C_d^+n_TaL_a$ and
$C_d^+m(B_0+N_c)x^{C_{\rm abs}\epsilon}$ by GFI1.
This is precisely the ledger's divisibility inequality.

For the place-error comparison the source gives, for each lift,
\[
 |E_{h,u,j}-A_T\log(g_0g_1)|\le C_d^+B_0,
 \qquad A_T=\sum_{t\in T_n}a_v(t).
\]
Sum over the $m$ lifts and apply the triangle inequality.  The
result is
\[
 \left|\sum_{j=0}^{m-1}E_{h,u,j}
                -mA_T\log(g_0g_1)\right|\le C_d^+mB_0,
                                                               \tag{GFI.16}
\]
the full \texttt{EComp5}.  Finally the source compares
$h_{h,u}+A_T\log T_{h,u}$ for any two words, with upper bound
$C_d^+(n_TaL_a+mB_0)$.  Pulling the word-dependent logarithm
inside the finite top sum gives the literal \texttt{HComp}.
% END GFI7

% BEGIN GFI8
\paragraph{GFI8: counting and both slope fields retain their normalization.}
\label{claim:GFI8}
Apply the supplied \texttt{Count\_at} to $D_{\mathrm{pt}(h)}$.
Replace its lift sum by $j=0,\ldots,m-1$, and its place sum by
$e_P$ using (GFI.5).  The factor $1/f$ remains unchanged; GFI1
identifies $m,k,v,L,L_a$ and $\log(g_0g_1)$.  The three resulting
inequalities are exactly \texttt{Count5}: the lower bound for
$\log T_{h,u}$, the upper bound $C_cmL$, and the cross-word
variation $C_c\log k$.

The two \texttt{SlopesM\_at} inequalities become
\[
 -C_sn_TaL_a(1+m/k)\le\mu_{M,h,u},\qquad
 \mu^{\max}_{M,h,u}\le C_sn_TaL_a,                    \tag{GFI.17}
\]
by (GFI.1) and GFI1.  They are the ledger's \texttt{Slopes}.
For the quotient slopes, reindexing the lift sum gives equality
of the original $\mathrm{Ee}$ with $\sum_jE_{h,u,j}$.
The upper and lower quotient bounds then have the same correction
$\mathrm{Ga}_h+\mathrm{Fo}_h-\mathrm{Ee}_{h,u}$ and the same
lower loss
\[
 C_1\mathrm{IO}_h+C_1n_Tm(a/k+b')\log T_{h,u}+C_1X''+1.
                                                               \tag{GFI.18}
\]
This supplies \texttt{QuotientSlopeBoundsWith}, with the original
constant $1$, the original $X''$, and the original lattice metrics.
% END GFI8

% BEGIN GFI9
\paragraph{GFI9: the loss witness supplies both ledger inequalities.}
\label{claim:GFI9}
Every $\mathrm{pt}(h)$ carries the pair in GFI2, so it satisfies
the same \texttt{Loss\_at} with constant $C_\ell$.  Its first
clause is, by GFI1 and (GFI.8),
\[
 \mathrm{IO}_h\le C_\ell n_Tmb\,x^{e_L\epsilon}.
\]
Its second clause quantifies over every permitted reading and
cutoff.  Apply it to the actual $A,T,a,w$ of GFI3 and the cutoff
$\lceil2a\rceil$; GFI5 checked every hypothesis of that application.
It gives
\[
 \mathrm{ASK}_h-\mathrm{Ga}_h
                  \le C_\ell n_Tmb\,x^{e_L\epsilon}.             \tag{GFI.19}
\]
The right-hand side is exactly $\mathscr D.\mathrm{LoP}(C)$.
These are both fields of \texttt{LossBounds}, simultaneously
for every $h$.  In particular no selected reading has been substituted
for the universally quantified statement before its hypotheses were
checked.
% END GFI9

% BEGIN GFI10
\paragraph{GFI10: frame, height, and defect on the same rational map.}
\label{claim:GFI10}
The definitions (GFI.7) show that the ledger's weighted mean $s_h$
is the source $s_{\rm Frame}$ and its excess $u_h$ is the source
$u_{\rm Frame}$.  The finite rank sum in the denominator is unchanged.
Thus the full all-height estimate in \texttt{FrameF}, specialized
to $HH_h=x^{C_H}b$, is precisely \texttt{SlopeIneq} whenever its
height hypotheses hold.  Those hypotheses do hold: \texttt{LinkDIV}
gives $s_h+\mu_{\max}(RR)\le x^{C_H}b$, and \texttt{LatF}
gives $\mu_{\min}(RR^\vee)=-\mu_{\max}(RR)$.
Hence
\[
 s_h-\mu^{\min}_{C,h}\le HH_h,
 \qquad 0<HH_h=x^{C_H}b.                               \tag{GFI.20}
\]
The upper height bound in \texttt{HHBound} is equality because
the record has $c_{HH}=C_H$.

The source's \texttt{LinkIRR} supplies
$0\le\delta_h\le\bar\delta$ for these actual retained rows.
The choices made after (GFI.14) supply
$0\le\eta_h\le\bar\eta$ and $0\le e_h\le\bar e$.
These six inequalities are \texttt{FParams}.  The frame estimates
and the height bound concern the same $RR^\vee$ and the same
original rational maps; no further dual or normalization is introduced.
% END GFI10

% BEGIN GFI11
\paragraph{GFI11: every error-window field and the trace-error budget.}
\label{claim:GFI11}
The chosen top density implies $n_T\ge1$: its lower bound is
$x^{-\epsilon/8}n/2>0$, and $n_T$ is an integer.  We have $x\ge1$,
$0<\epsilon\le1$, $b\ge1$, $a=xb$, and $k\ge1$.
The original windows of GFA11 give
\[
\begin{gathered}
 m\le x^{1/2+\epsilon/2},\quad xm\le k\le x^{e_K},\quad
 k\le a,\quad 0\le L_a\le2x^\epsilon,\\
 kb'\le a,\quad X''\le a,\quad
 \log\#W_h\le(k-3)\log2,
 \qquad e_K\ge0.                                      \tag{GFI.21}
\end{gathered}
\]
Also $0\le L=x^\epsilon\le2x^\epsilon$, $b'>0$, $X''\ge0$,
and $B_0,N_c\ge0$.  Top nonnegativity and the upper bound $2a$
are (GFI.10).  Since $x^{\epsilon/(7J)}\ge1$,
$\bar\delta\le e^{-1}\le1/2$.  Every nonnegative coefficient
in (GFI.13) is at most $K$, and the three exponent parameters
are nonnegative.  Together with (GFI.14) these prove every field
of \texttt{ErrWindows0}; in particular the word-cardinality bound
is at every actual member.

For the extra field of \texttt{ErrWindows}, (GFI.14) gives
\[
 (\bar\eta+4\sqrt{\bar e})x^{C_H}b\le6b
 \le Kx^\epsilon n_Ta.                                \tag{GFI.22}
\]
The second inequality uses only $K\ge6$, $x^\epsilon\ge1$,
$n_T\ge1$, and $a=xb\ge b$.  It proves the full trace-error
budget with the very same $K$ and $C_H$.  No weaker window
statement is substituted at this step.
% END GFI11

% BEGIN GFI12
\paragraph{GFI12: the remaining growth, cost, and size records.}
\label{claim:GFI12}
The actual definitions of the two budgets give equalities
\[
 B_0=2nmbx^\epsilon,\qquad
 N_c=\frac{\#T_2\,kb x^\epsilon}{n(nx^{-\epsilon})}.      \tag{GFI.23}
\]
For the first use $f=2n$; for the second use (GFI.4).
Thus the two \texttt{BNBounds} inequalities hold with equality.

Put $\eta=\eta_{\rm Sz}$, $X'=e^{x^{\epsilon/3}}$ and
$\theta=e^{-c_6x^{\epsilon/8}}$.  These have the required signs,
and $b'X'=a=xb$.  The near-$x$ top condition is (GFI.10).
The node conditions follow from GFA10 after transporting the same
second set and the same $\eta_{\rm Sz}$ using GFI1.  GFI5 supplies
the common cutoff condition.  These are all fields of
\texttt{GrowthInequalitySizes}.  Applying the last two clauses
of the single character record from GFI5 to the supplied
\texttt{ChosenThin} and \texttt{ChosenQD} gives
\texttt{ThinScales} and \texttt{ValuesAtX} with this $\theta$
and $C_L=C'_5$.  Therefore the full remaining growth record holds.

The same windows give
\[
 2f(x+1)\theta\le1/x,\qquad
 6f(x+1)(\eta+1)\le b/x,                                \tag{GFI.24}
\]
and the place cardinality gives $f_{\rm Pl}(H)=f$, proving
\texttt{CostWindows}.  Finally the five \texttt{RestSize}
conditions are exactly
\[
 a\le2xb,\quad x^{-\epsilon/8}n/2\le n_T,\quad
 x^{-\epsilon}\le\#T_2/n,\quad
 x^{1/2-\epsilon}\le m,\quad k\le n^2x^{1/2-\epsilon}.
                                                               \tag{GFI.25}
\]
They follow, respectively, from $a=xb$, the original top density,
the second-set window and (GFI.4), and the final two rank windows.
All these conclusions concern the same ledger and cutoff.
% END GFI12

% BEGIN GFI13
\paragraph{GFI13: the literal \texttt{ChainTied} exists eventually.}
\label{claim:GFI13}
Under the hypotheses (GFI.11) and GFI6, there are constants
$K,c_{HH},e_K,C_L$ and a single threshold $x_0$ such that
\[
 x\ge x_0,\quad x\in\mathcal F.\mathrm{Win}
 \quad\Longrightarrow\quad
 \texttt{ChainTied}(\mathcal F,C_{\rm abs},e_L,
                              K,c_{HH},e_K,C_L,x).               \tag{GFI.26}
\]
Indeed choose the constants and threshold as in GFI6.  For each such
$x$ use the common model of GFI2.  Its original frame, rank and
word sets, top type $T_n$, constants (GFI.13), ledger (GFI.7)--(GFI.9),
sets (GFI.3), equivalences $e_T,e_P$, character functional and cutoff
from GFI5 are the witnesses in \texttt{ChainTied}.

The four tie predicates are GFI4--GFI5.  The five scalar/constant
identities and inequalities are definitions: $\mathscr D.x=x$,
$\mathscr D.\epsilon=\epsilon$, $e_D=C_{\rm abs}$,
$e_P=e_L$ and $c_{HH}=C_H$.  The \texttt{Ident} record is GFI5.
The eleven ledger fields, in their source order, are GFI6;
the three results of GFI7; the three results of GFI8; GFI9;
and the three results of GFI10.  GFI11 supplies
\texttt{ErrWindows}; GFI12 supplies \texttt{BNBounds}, the
remaining growth hypotheses, \texttt{CostWindows}, and
\texttt{RestSize}.  This checks every field of the original
existential record and proves (GFI.26).

In particular, the constants are chosen before $x$, although the
realized degree pair, base index and error bounds may depend on $x$
as the source permits.  A scalar bound established only after
varying the base point would not give this conclusion.
% END GFI13

% BEGIN GFI14
\paragraph{GFI14: the original chosen inputs supply the small-share inputs.}
\label{claim:GFI14}
Suppose $\texttt{ChosenE}(\mathcal F,C_0,C_H)$ and the same
family's \texttt{ChosenWall} are supplied.  Unpack the original
positive constants $C'_5,c_6,C_b$.  The top, quadratic, thin,
proviso and code predicates are already present.  We verify the
three other fields of \texttt{WeightVectorSmallSharesSeven}.

First the original second-set lower bound is
\[
 \#T_2^{\rm raw}\ge
 \frac{n_0}{L_w}
 \exp\!\left[-C_{T2}\sqrt{\log n_0\log L_w}\right],
 \qquad n=o n_0.                                      \tag{GFI.27}
\]
The family fixes $C_{T2}$, bounds $o\le C_o$, and has
$1\le n_0\le n\le x^{9/4}$ and $2\le L_w\le(\log x)^2$
eventually, by its original size record.  Set
$C=\max(C_{T2},0)$ and $O=\max(C_o,1)$.  For $x\ge e$,
the logarithm of the inverse lower density in (GFI.27) is at most
\[
 \log O+2\log\log x+
       C\sqrt{\tfrac92\log x\log\log x}.              \tag{GFI.28}
\]
Choose a single threshold making each of the three terms at most
$\epsilon\log x/3$.  This is possible because
$1/\log x\to0$ and $\log\log x/\log x\to0$; for the square-root
term the required condition is
$(9/2)C^2\log\log x/\log x\le\epsilon^2/9$.
Dividing (GFI.27) by $n$ and exponentiating gives
$x^{-\epsilon}\le\#T_2^{\rm raw}/n$ uniformly.  This is exactly
\texttt{ThickSetDensityEventually}.

Next the source \texttt{ChosenBad} gives its original bad-top sum
at most
$4C_bx^{-2}\delta_1n\sum_{p\in\mathcal P}\log p/p$, where
$\mathcal P$ is the original large-order residue-prime set.
Since $\delta_1n\le2n_T$, it suffices to bound that prime sum.
An elementary bound is enough.  For $X\ge e$, putting
$M=\lceil X\rceil$ and using $\sum_{r=1}^M1/r\le1+\log M$ gives
\[
 \sum_{\substack{p<M\\p\ {
m prime}}}\frac{\log p}{p}
 \le \log M(1+\log M)\le8(\log X)^2.                  \tag{GFI.29}
\]
Here $M\le2X$, so $\log M\le2\log X$ and
$1+\log M\le3\log X$.  Apply (GFI.29) to $X=X'$.
The bad-top sum is at most
$64C_bn_Tx^{-2+2\epsilon/3}$.  After imposing
$64C_b\le x^{\epsilon/3}$, this is at most
$n_Tx^{-2+\epsilon}$, the literal \texttt{ChosenBadp} bound.
No distribution theorem for primes is needed for this conversion.

Finally, at a late allowed scale choose an existing index $i_0$.
Its actual degrees define integer functions on $T_n$, extended by
zero elsewhere; injectivity of the labels makes them well-defined.
The original top window and $b_T\le b_{\rm top}$ give admissibility.
The same $i_0$ is a realizing index.  Every other carrier satisfies
\texttt{ChW\_at} by \texttt{ChosenWall}.  This is
\texttt{ChosenWS}, also as in HW3.  Combining these fields with the
unchanged five chosen inputs proves the entire original Seven record.
Its constants and thresholds retain the required order.
% END GFI14

% BEGIN GFI15
\paragraph{GFI15: from the eleven original family fields to \texttt{LinkGLOBAL}.}
\label{claim:GFI15}
Fix $C_{\rm abs}\ge0$, $C_H>0$, $C_0\ge3$, and suppose the
eleven original \texttt{FamGLmin} fields with upper-column constant
$4$ hold for this family.  Supply \texttt{TuSize},
$\texttt{WeightVectorSmallSharesSeven}(\mathcal F,C_0,C_H)$ and
\[
 \texttt{LinkG},\quad\texttt{LinkS}(C_0),\quad
 \texttt{LinkIRR},\quad\texttt{LinkDIV}(C_{\rm abs},C_H),\quad
 \texttt{LinkMu}.                                       \tag{GFI.30}
\]
Then the literal \texttt{LinkGLOBAL7} conclusion holds with
$e_D=C_{\rm abs}$ and $e_L=0$.

Here is the full dependency installation.  The original guarded
local fields are activated by \texttt{TuSize} through GLE2.
We first recover the three ramified-model conclusions directly
from the supplied $h_W,h_\mu$; the stronger second-layer hypotheses
of the RSL construction are not required for this generic consumer.
Fix an original word, put $T=(D.\mathrm{mem}(u)).A.T$, and take
its supplied \texttt{FreeT} record.  Its finite set $X$ has a free
action of
\[
 \mathscr G=\prod_{v\mid T}k_v^\times,
 \qquad |k_v|=T,\qquad |\{v\mid T\}|=f.
\]
Choose one representative of each orbit.  If there are $r$ orbits,
then $X$ is a disjoint union of $r$ regular $\mathscr G$-sets.
The stored equivariant linear equivalence from the original full
form space to $\mathbb C^X$ carries each original norm operator
to summation in its coordinate.  On one orbit the function space
is $\bigotimes_{v\mid T}\mathbb C^{T-1}$.  Split each factor as
the kernel of coordinate summation plus the line spanned by one
coordinate vector of sum one.  Expanding the tensor product shows
that the joint kernel is the tensor product of those kernels,
whose dimension is $(T-2)^f$.  Thus, with the original dimensions,
\[
 D.duS(u)=r(T-1)^f,\qquad
 d_u=(D.\mathrm{mem}(u)).A.du=r(T-2)^f,
 \qquad d_u=D.duS(u)\left(1-\frac1{T-1}\right)^f.
                                                               \tag{GFI.30a}
\]
The identifications with these dimensions use the original
\texttt{emb\_basis}, the actual joint-kernel definition of the
ramified forms, and the range of the original $D.\mathrm{nuIn}$.
No integral or isometric version of the free model is needed for
this dimension computation.

The supplied $h_\mu$ is already the exact mean identity
\[
 \mu(\mathrm{ramLat})=\mu(\mathrm{lat}(0))
                         -\frac{f\log(T-1)}{T-2}.
                                                               \tag{GFI.30b}
\]
The original \texttt{ChosenProviso}, contained in the Seven packet,
forces $D.duS(u)>0$: if it were zero, its normalized lower bound
would read $C_0^{-1}\le0$, whereas $C_0>0$.
The setting gives $T>2fN(m)+1>2$ since $f,N(m)\ge1$.
Consequently (GFI.30a) gives $d_u>0$.
The actual ramified lattice is the intersection of the full original
lattice with its rational ramified subspace, with restricted metrics.
Every positive-rank lattice used in its maximal-slope supremum is
therefore admissible for the full lattice's supremum.  If saturation
is part of the convention, this remains true because the ramified
intersection is saturated and saturation inside it is saturated in
the full lattice.  This proves
$\mu_{\max}(\mathrm{ramLat})\le\mu_{\max}(\mathrm{lat}(0))$
by the same inclusion argument written in RSL10.  Together with
(GFI.30a)--(GFI.30b), these are all three literal
\texttt{family\_property\_and} clauses.

Next $h_I,h_R,h_B$ give the original integral and ramified operator
identities and second-branch inclusion/length.  The finite-DVR and
ordered-product arguments of GLE10--GLE11 consume exactly these
clauses and the setting, giving the full original congruence length
bound and $N_\lambda\subset S_{\rm ram}$ at every place over $T$.
RSL11 then gives the original commensurability and fullness, using
the already stored \texttt{nuDef} and \texttt{nprimeDef}.
Their ranks equal $d_u$ by the original rational embedding, so
(GFI.30a) and the proviso give \texttt{J.MuRank} as well as the
ramified rank identity.  All thresholds used in this paragraph are
maxima of the supplied family thresholds and precede every word
and place.  In particular no additional second-layer selector
has entered this implication.

GVC1--GVC18 use the mass, volume, numerical and column fields to
give \texttt{CountF} and \texttt{SlopesMF}; their derivative
window is the existing family-size projection.  This direct route
does not assume the stronger unused \texttt{PartFiveFamilyBounds}.
GFA1--GFA5, using (GFI.30), give \texttt{SlopesNF}.
GFA6 gives the positive target rank and \texttt{LatF}.
GFA7--GFA9 use the same Bessel map, proviso, code, local ranks,
and links to give the full \texttt{FrameF}$(C_H)$.
GFA10--GFA12 give \texttt{NodesF}, \texttt{WinF} and
\texttt{CharF}, using the supplied thick density for the window.
LFS17, with the supplied \texttt{ChosenTT},
\texttt{ChosenBadp}, \texttt{ChosenWS} and \texttt{LinkG}, gives
the original \texttt{LossFS}$(0)$.

These are exactly (GFI.11).  GFI13, with the unchanged chosen
quadratic and thin constants $C'_5,c_6$, now proves
\texttt{LinkGLOBAL7}.  Notice that $C'_5=C_L$ is the quadratic
constant; it need not equal the upper-column constant $4$.

If instead \texttt{ChosenE} and \texttt{ChosenWall} are supplied,
GFI14 constructs the required original Seven record on this same
family.  Applying the preceding argument to every such
\texttt{ChosenE} and every (GFI.30) proves the literal
\texttt{LinkGLOBAL}$(\mathcal F,C_{\rm abs},C_H,C_0,C_{\rm abs},0)$.
The retained \texttt{DGamSize} premise in the final exported
signature is allowed; it is used upstream when MT6 constructs
\texttt{LinkMu}, and is not consumed a second time here.
% END GFI15

% BEGIN GFI16
\paragraph{GFI16: the \texttt{Closed} quantifiers are preserved.}
\label{claim:GFI16}
Let $\mathcal C$ be the same original family class.  Suppose
$\texttt{FamGLmin}(\mathcal C,4)$ holds and every family in
$\mathcal C$ has its original \texttt{ChosenWall}.  Then
\[
 \texttt{Closed}_{\mathcal C}\left(
 \texttt{TuSize}\Longrightarrow\texttt{DGamSize}
 \Longrightarrow
 \texttt{LinkGLOBAL}(C_{\rm abs},C_H,C_0,C_{\rm abs},0)
 \right)                                                \tag{GFI.31}
\]
for $C_{\rm abs}\ge0$ and $C_0\ge3$.

To check the exact order, write each of the eleven \texttt{Closed}
inputs as a positive cutoff $\epsilon_j$ followed, for every
permitted $\epsilon$, by a transported tuple $\mathcal K_j$.
Take the minimum of these finitely many cutoffs and a fixed
positive cutoff below $1/52$.  For that $\epsilon$ take the
componentwise supremum of the tuples.  Its \texttt{pre} component
is a maximum, and its \texttt{t}, \texttt{mid} and \texttt{last}
components are pointwise maxima of functions of the original lists.
Thus no choice here depends on the later family or on $C_H$.
The definition of \texttt{Large} and monotonicity recover each
original largeness input for the same family and lists.

Join also the tuple with \texttt{pre}$=0$, \texttt{t}$=0$,
\texttt{mid}$=1$ and \texttt{last}$=0$.  Because $C_H$ is one
of the original middle constants, its largeness condition forces
$C_H\ge1$.  This proves the strictly positive height exponent
needed in GFI15 before any frame estimate is used.

Now take any supplied $\mathcal F\in\mathcal C$, $C_H$ satisfying
that joined largeness condition.  All eleven fields and
\texttt{ChosenWall} concern this same family.  Apply GFI15 to
any supplied \texttt{TuSize}, \texttt{DGamSize}, chosen data and
global links.  Its eventual $x$-threshold may depend on those
objects, exactly as \texttt{LinkGLOBAL} permits, but precedes
every index and word.  This is (GFI.31) with the literal
\texttt{Closed} quantifier order.
% END GFI16

% BEGIN GFI17
\paragraph{GFI17: completion of the selected global assembly obligation.}
\label{claim:GFI17}
Retain the original class $\mathcal A$, its \texttt{Duty} on every
member, and the compatible second-layer assertion on
$\mathcal C=D_3(3,\mathcal A)$ from GLS13.  That theorem gives
the full eleven-field $\texttt{FamGLmin}(\mathcal C,4)$.
Membership in $D_3(3,\mathcal A)$ projects to membership in
$\mathcal A$, and the original Duty record projects to
\texttt{ChosenWall} for that same family.  GFI16, with
$C_{\rm abs}=2$ and $C_0=3$, therefore proves
\[
 \boxed{\quad
 \texttt{Closed}_{D_3(3,\mathcal A)}
  \bigl(\texttt{TuSize}\to\texttt{DGamSize}\to
          \texttt{LinkGLOBAL}(2,C_H,3,2,0)\bigr).
 \quad}                                                   \tag{GFI.32}
\]
This is the full remaining numerical/lattice consumer in
\texttt{TASK\_GL\_FINAL\_ASSEMBLY}.  All its loss and ledger
conditions now have human proofs, and $e_D=2,e_L=0$ are unchanged.
When the final common-family inputs are supplied, ZW1--ZW4 apply
to this very \texttt{ChainTied} and retain the final exponent
$1/2+9\epsilon$.  No numerical absorption needs to be added
between (GFI.32) and that existing endpoint.

The conclusion is conditional on the same common construction,
compatible selector and duties as GLS13.  It supplies none of
their existence, nor the actual denominator, retained-defect,
period-height, local-gain or spectral producers.  Those remain
explicit in the final dependency audit.  In particular (GFI.32)
closes the assembly obligation without asserting a complete human
proof of either final theorem.
% END GFI17
